\documentclass[12pt]{article}
\usepackage[T1]{fontenc}
\usepackage[utf8]{inputenc}
\usepackage{lmodern}
\usepackage{textcomp}
\usepackage{enumitem}
\usepackage{geometry}
\geometry{margin=1in}
\begin{document}
\begin{center}
Answer Key - Mathematics - Graduate Level Assessment 1 \
\small (Answers are ordered exactly as in the question paper.)
\end{center}

\section*{Multiple Choice}
\begin{enumerate}[label=\arabic*.]
\item \textbf{Answer:} A
\\ \textit{Options:} A) a circle; B) an infinite interval; C) the negative real axis; D) the positive real axis; E) a ray
\\ \textit{Note:} Correct option: A - a circle
\item \textbf{Answer:} A
\\ \textit{Options:} A) \$125.00; B) \$20.70; C) \$130.00; D) \$133.00; E) \$115.00
\\ \textit{Note:} Correct option: A - \$125.00
\item \textbf{Answer:} B
\\ \textit{Options:} A) 11.234567; B) 10.065778; C) 6.789012; D) 7.654321; E) 9.876543
\\ \textit{Note:} Correct option: B - 10.065778
\item \textbf{Answer:} A
\\ \textit{Options:} A) 16 over 36; B) 9 over 4; C) 8 over 18; D) 48 over 108; E) 24 over 54
\\ \textit{Note:} Correct option: A - 16 over 36
\item \textbf{Answer:} A
\\ \textit{Options:} A) 9; B) 7; C) 8; D) 1; E) 4
\\ \textit{Note:} Correct option: A - 9
\end{enumerate}

\section*{True / False}
\begin{enumerate}[label=\arabic*.]
\setcounter{enumi}{5}
\item \textbf{Answer:} False
\\ \textit{Options:} A) True; B) False
\\ \textit{Note:} LLM-generated false statement. Correct answer: '24'.
\item \textbf{Answer:} True
\\ \textit{Options:} A) True; B) False
\\ \textit{Note:} LLM-generated true statement based on MCQ.
\item \textbf{Answer:} False
\\ \textit{Options:} A) True; B) False
\\ \textit{Note:} LLM-generated false statement. Correct answer: '0.4295'.
\item \textbf{Answer:} False
\\ \textit{Options:} A) True; B) False
\\ \textit{Note:} LLM-generated false statement. Correct answer: '56'.
\item \textbf{Answer:} True
\\ \textit{Options:} A) True; B) False
\\ \textit{Note:} LLM-generated true statement based on MCQ.
\end{enumerate}

\section*{Long Form Response}
\begin{enumerate}[label=\arabic*.]
\setcounter{enumi}{10}
\item \textbf{Answer:} If two numbers give the same remainder when divided by 5, they are said to be equivalent, modulo 5. From $n^2$ to $n^3$, we have multiplied by $n$. Since $n^2$ is equivalent to 4 (modulo 5), and $n^3$ is equivalent to 2 (modulo 5), we are looking for an integer $n$ for which $4\cdot n$ is equivalent to 2, modulo 5. Notice that if $n$ is greater than 4, then we can replace it with its remainder when divided by 5 without changing whether it satisfies the condition. Therefore, we may assume that $0\leq n <5$. Trying 0, 1, 2, 3, and 4, we find that only $\boxed{3}$ times 4 leaves a remainder of 2 when divided by 5.
\\ \textit{Note:} Students should show all steps in their mathematical solution.
\item \textbf{Answer:} 16. Explain why this is the correct answer and provide supporting evidence. Show the calculation or reasoning that leads to this conclusion.
\\ \textit{Note:} Long-form derived from MCQ; award credit for stating the correct idea and giving supporting reasoning.
\end{enumerate}

\vfill
\noindent\textit{End of Answer Key}
\end{document}